\documentclass[12pt]{article}

\usepackage{amsmath}
\usepackage{amssymb}
\usepackage[margin=1in]{geometry}

\title{APM1220: Applied Multivariate Data Analysis\\FA 1}
\author{}
\date{}

\begin{document}

\maketitle

\section*{Problem 2.2}

Given
\[
A =
\begin{bmatrix}
-1 & 3 \\
4 & 2
\end{bmatrix},
\qquad
B =
\begin{bmatrix}
4 & -3 \\
1 & -2 \\
-2 & 0
\end{bmatrix},
\qquad
C =
\begin{bmatrix}
5 \\
-4 \\
2
\end{bmatrix}.
\]

\subsection*{(a) $5A$}

\[
5A
=
5
\begin{bmatrix}
-1 & 3 \\
4 & 2
\end{bmatrix}
=
\begin{bmatrix}
-5 & 15 \\
20 & 10
\end{bmatrix}.
\]

Therefore,
\[
\boxed{
5A =
\begin{bmatrix}
-5 & 15 \\
20 & 10
\end{bmatrix}
}
\]

\subsection*{(b) $BA$}

The dimensions are
\[
B_{3\times2}A_{2\times2},
\]
so the multiplication is defined and the result is a $3\times2$ matrix.

\[
BA
=
\begin{bmatrix}
4 & -3 \\
1 & -2 \\
-2 & 0
\end{bmatrix}
\begin{bmatrix}
-1 & 3 \\
4 & 2
\end{bmatrix}.
\]

\[
BA
=
\begin{bmatrix}
4(-1)+(-3)(4) & 4(3)+(-3)(2) \\
1(-1)+(-2)(4) & 1(3)+(-2)(2) \\
(-2)(-1)+0(4) & (-2)(3)+0(2)
\end{bmatrix}
\]

\[
=
\begin{bmatrix}
-16 & 6 \\
-9 & -1 \\
2 & -6
\end{bmatrix}.
\]

Therefore,
\[
\boxed{
BA =
\begin{bmatrix}
-16 & 6 \\
-9 & -1 \\
2 & -6
\end{bmatrix}
}
\]

\subsection*{(c) $A'B'$}

First, take the transpose of each matrix:

\[
A'
=
\begin{bmatrix}
-1 & 4 \\
3 & 2
\end{bmatrix},
\qquad
B'
=
\begin{bmatrix}
4 & 1 & -2 \\
-3 & -2 & 0
\end{bmatrix}.
\]

Therefore,

\[
A'B'
=
\begin{bmatrix}
-1 & 4 \\
3 & 2
\end{bmatrix}
\begin{bmatrix}
4 & 1 & -2 \\
-3 & -2 & 0
\end{bmatrix}.
\]

\[
A'B'
=
\begin{bmatrix}
(-1)(4)+4(-3) &
(-1)(1)+4(-2) &
(-1)(-2)+4(0) \\
3(4)+2(-3) &
3(1)+2(-2) &
3(-2)+2(0)
\end{bmatrix}
\]

\[
=
\begin{bmatrix}
-16 & -9 & 2 \\
6 & -1 & -6
\end{bmatrix}.
\]

Therefore,

\[
\boxed{
A'B'
=
\begin{bmatrix}
-16 & -9 & 2 \\
6 & -1 & -6
\end{bmatrix}
}
\]

\subsection*{(d) $C'B$}

First,

\[
C'
=
\begin{bmatrix}
5 & -4 & 2
\end{bmatrix}.
\]

Therefore,

\[
C'B
=
\begin{bmatrix}
5 & -4 & 2
\end{bmatrix}
\begin{bmatrix}
4 & -3 \\
1 & -2 \\
-2 & 0
\end{bmatrix}.
\]

\[
C'B
=
\begin{bmatrix}
5(4)+(-4)(1)+2(-2) &
5(-3)+(-4)(-2)+2(0)
\end{bmatrix}.
\]

\[
=
\begin{bmatrix}
12 & -7
\end{bmatrix}.
\]

Therefore,

\[
\boxed{
C'B =
\begin{bmatrix}
12 & -7
\end{bmatrix}
}
\]

\subsection*{(e) Is $AB$ defined?}

The dimensions are

\[
A_{2\times2}
\qquad\text{and}\qquad
B_{3\times2}.
\]

For $AB$ to be defined, the number of columns of $A$ must equal the number of rows of $B$.

Here,

\[
2 \neq 3.
\]

Therefore,

\[
\boxed{AB\text{ is not defined}.}
\]


\section*{Problem 2.4}

Suppose $A^{-1}$ and $B^{-1}$ exist.

\subsection*{(a) Prove that $(A')^{-1}=(A^{-1})'$}

Since $A^{-1}$ exists,

\[
AA^{-1}=I.
\]

Taking the transpose of both sides gives

\[
(AA^{-1})'=I'.
\]

Using the properties

\[
(XY)'=Y'X'
\]

and

\[
I'=I,
\]

we obtain

\[
(A^{-1})'A'=I.
\]

Similarly, since

\[
A^{-1}A=I,
\]

taking the transpose gives

\[
(A^{-1}A)'=I',
\]

so

\[
A'(A^{-1})'=I.
\]

Therefore, $(A^{-1})'$ is both a left and right inverse of $A'$. Hence,

\[
\boxed{(A')^{-1}=(A^{-1})'}.
\]

\subsection*{(b) Prove that $(AB)^{-1}=B^{-1}A^{-1}$}

Consider

\[
(B^{-1}A^{-1})(AB).
\]

Using associativity,

\[
(B^{-1}A^{-1})(AB)
=
B^{-1}(A^{-1}A)B.
\]

Since

\[
A^{-1}A=I,
\]

we have

\[
B^{-1}(A^{-1}A)B
=
B^{-1}IB
=
B^{-1}B
=
I.
\]

We also check the other order:

\[
(AB)(B^{-1}A^{-1})
=
A(BB^{-1})A^{-1}.
\]

Since

\[
BB^{-1}=I,
\]

we obtain

\[
AIA^{-1}
=
AA^{-1}
=
I.
\]

Thus, $B^{-1}A^{-1}$ is the inverse of $AB$. Therefore,

\[
\boxed{(AB)^{-1}=B^{-1}A^{-1}}.
\]


\section*{Problem 2.8}

Given

\[
A=
\begin{bmatrix}
1 & 2\\
2 & -2
\end{bmatrix},
\]

find the eigenvalues $\lambda_1$ and $\lambda_2$, the associated normalized eigenvectors $e_1$ and $e_2$, and the spectral decomposition of $A$.

\subsection*{Step 1: Find the eigenvalues}

The eigenvalues satisfy

\[
\det(A-\lambda I)=0.
\]

We have

\[
A-\lambda I
=
\begin{bmatrix}
1-\lambda & 2\\
2 & -2-\lambda
\end{bmatrix}.
\]

Therefore,

\[
\det(A-\lambda I)
=
(1-\lambda)(-2-\lambda)-4.
\]

Expanding,

\[
(1-\lambda)(-2-\lambda)-4
=
\lambda^2+\lambda-6.
\]

Thus,

\[
\lambda^2+\lambda-6=0.
\]

Factoring,

\[
(\lambda-2)(\lambda+3)=0.
\]

Therefore,

\[
\boxed{\lambda_1=2,\qquad \lambda_2=-3}.
\]

\subsection*{Step 2: Find the eigenvector associated with $\lambda_1=2$}

Solve

\[
(A-2I)e_1=0.
\]

We have

\[
A-2I
=
\begin{bmatrix}
-1 & 2\\
2 & -4
\end{bmatrix}.
\]

Let

\[
e_1=
\begin{bmatrix}
x\\
y
\end{bmatrix}.
\]

Then

\[
-x+2y=0,
\]

so

\[
x=2y.
\]

Choose $y=1$. Then $x=2$, giving the eigenvector

\[
\begin{bmatrix}
2\\
1
\end{bmatrix}.
\]

Its norm is

\[
\sqrt{2^2+1^2}
=
\sqrt{5}.
\]

Therefore, the normalized eigenvector is

\[
\boxed{
e_1=
\frac{1}{\sqrt{5}}
\begin{bmatrix}
2\\
1
\end{bmatrix}
}.
\]

\subsection*{Step 3: Find the eigenvector associated with $\lambda_2=-3$}

Solve

\[
(A+3I)e_2=0.
\]

We have

\[
A+3I
=
\begin{bmatrix}
4 & 2\\
2 & 1
\end{bmatrix}.
\]

Let

\[
e_2=
\begin{bmatrix}
x\\
y
\end{bmatrix}.
\]

Then

\[
4x+2y=0,
\]

which gives

\[
y=-2x.
\]

Choose $x=1$. Then $y=-2$, giving

\[
\begin{bmatrix}
1\\
-2
\end{bmatrix}.
\]

Its norm is

\[
\sqrt{1^2+(-2)^2}
=
\sqrt{5}.
\]

Therefore,

\[
\boxed{
e_2=
\frac{1}{\sqrt{5}}
\begin{bmatrix}
1\\
-2
\end{bmatrix}
}.
\]

\subsection*{Step 4: Spectral decomposition}

For a symmetric matrix, the spectral decomposition is

\[
A=\lambda_1e_1e_1'+\lambda_2e_2e_2'.
\]

First,

\[
e_1e_1'
=
\frac{1}{5}
\begin{bmatrix}
2\\
1
\end{bmatrix}
\begin{bmatrix}
2&1
\end{bmatrix}
=
\frac{1}{5}
\begin{bmatrix}
4&2\\
2&1
\end{bmatrix}.
\]

Similarly,

\[
e_2e_2'
=
\frac{1}{5}
\begin{bmatrix}
1\\
-2
\end{bmatrix}
\begin{bmatrix}
1&-2
\end{bmatrix}
=
\frac{1}{5}
\begin{bmatrix}
1&-2\\
-2&4
\end{bmatrix}.
\]

Therefore,

\[
A
=
2
\left(
\frac{1}{5}
\begin{bmatrix}
4&2\\
2&1
\end{bmatrix}
\right)
-
3
\left(
\frac{1}{5}
\begin{bmatrix}
1&-2\\
-2&4
\end{bmatrix}
\right).
\]

Hence,

\[
\boxed{
A=2e_1e_1'-3e_2e_2'
}.
\]


\section*{Problem 2.11}

Let $A=\{a_{ij}\}$ be a $p\times p$ diagonal matrix, where

\[
a_{ij}=0,\qquad i\neq j.
\]

Thus,

\[
A=
\begin{bmatrix}
a_{11}&0&\cdots&0\\
0&a_{22}&\cdots&0\\
\vdots&\vdots&\ddots&\vdots\\
0&0&\cdots&a_{pp}
\end{bmatrix}.
\]

Expanding the determinant along the first row,

\[
|A|
=
a_{11}A_{11}+0+\cdots+0
=
a_{11}A_{11}.
\]

The submatrix corresponding to $A_{11}$, obtained by deleting the first row and first column, is also diagonal:

\[
A_{11}
=
\begin{bmatrix}
a_{22}&0&\cdots&0\\
0&a_{33}&\cdots&0\\
\vdots&\vdots&\ddots&\vdots\\
0&0&\cdots&a_{pp}
\end{bmatrix}.
\]

Expanding its determinant along the first row gives

\[
A_{11}=a_{22}A_{22}.
\]

Therefore,

\[
|A|=a_{11}a_{22}A_{22}.
\]

Repeating this process for each successive submatrix gives

\[
|A|
=
a_{11}a_{22}a_{33}\cdots a_{pp}.
\]

Hence,

\[
\boxed{
|A|=a_{11}a_{22}\cdots a_{pp}
}.
\]


\section*{Problem 2.24}

Given

\[
\Sigma=
\begin{bmatrix}
4&0&0\\
0&9&0\\
0&0&1
\end{bmatrix}.
\]

\subsection*{(a) Find $\Sigma^{-1}$}

Since $\Sigma$ is diagonal, its inverse is obtained by taking the reciprocal of each nonzero diagonal element:

\[
\boxed{
\Sigma^{-1}
=
\begin{bmatrix}
\frac14&0&0\\
0&\frac19&0\\
0&0&1
\end{bmatrix}
}.
\]

\subsection*{(b) Eigenvalues and eigenvectors of $\Sigma$}

Because $\Sigma$ is diagonal, its eigenvalues are its diagonal elements:

\[
\boxed{
\lambda_1=4,\qquad
\lambda_2=9,\qquad
\lambda_3=1
}.
\]

For $\lambda_1=4$,

\[
\boxed{
e_1=
\begin{bmatrix}
1\\
0\\
0
\end{bmatrix}
}.
\]

For $\lambda_2=9$,

\[
\boxed{
e_2=
\begin{bmatrix}
0\\
1\\
0
\end{bmatrix}
}.
\]

For $\lambda_3=1$,

\[
\boxed{
e_3=
\begin{bmatrix}
0\\
0\\
1
\end{bmatrix}
}.
\]

Thus,

\[
\boxed{
(\lambda_1,\lambda_2,\lambda_3)=(4,9,1)
}
\]

with corresponding eigenvectors

\[
\boxed{
e_1=
\begin{bmatrix}
1\\
0\\
0
\end{bmatrix},
\quad
e_2=
\begin{bmatrix}
0\\
1\\
0
\end{bmatrix},
\quad
e_3=
\begin{bmatrix}
0\\
0\\
1
\end{bmatrix}.
}
\]

\subsection*{(c) Eigenvalues and eigenvectors of $\Sigma^{-1}$}

We have

\[
\Sigma^{-1}
=
\begin{bmatrix}
\frac14&0&0\\
0&\frac19&0\\
0&0&1
\end{bmatrix}.
\]

Therefore, its eigenvalues are

\[
\boxed{
\lambda_1=\frac14,\qquad
\lambda_2=\frac19,\qquad
\lambda_3=1
}.
\]

The corresponding eigenvectors remain the same:

\[
\boxed{
e_1=
\begin{bmatrix}
1\\
0\\
0
\end{bmatrix},
\quad
e_2=
\begin{bmatrix}
0\\
1\\
0
\end{bmatrix},
\quad
e_3=
\begin{bmatrix}
0\\
0\\
1
\end{bmatrix}.
}
\]


\section*{Problem 2.30}

Given

\[
X=
\begin{bmatrix}
X_1\\
X_2\\
X_3\\
X_4
\end{bmatrix},
\qquad
\mu_X=
\begin{bmatrix}
4\\
3\\
2\\
1
\end{bmatrix},
\]

and

\[
\Sigma_X=
\begin{bmatrix}
3&0&2&2\\
0&1&1&0\\
2&1&9&-2\\
2&0&-2&4
\end{bmatrix}.
\]

Partition $X$ as

\[
X=
\begin{bmatrix}
X^{(1)}\\
X^{(2)}
\end{bmatrix},
\]

where

\[
X^{(1)}
=
\begin{bmatrix}
X_1\\
X_2
\end{bmatrix},
\qquad
X^{(2)}
=
\begin{bmatrix}
X_3\\
X_4
\end{bmatrix}.
\]

Also,

\[
A=
\begin{bmatrix}
1&2
\end{bmatrix},
\qquad
B=
\begin{bmatrix}
1&-2\\
2&-1
\end{bmatrix}.
\]

\subsection*{(a) $E(X^{(1)})$}

The first two elements of the mean vector are

\[
\boxed{
E(X^{(1)})
=
\begin{bmatrix}
4\\
3
\end{bmatrix}
}.
\]

\subsection*{(b) $E(AX^{(1)})$}

Using

\[
E(AX^{(1)})=AE(X^{(1)}),
\]

we have

\[
E(AX^{(1)})
=
\begin{bmatrix}
1&2
\end{bmatrix}
\begin{bmatrix}
4\\
3
\end{bmatrix}.
\]

Therefore,

\[
E(AX^{(1)})=4+6=10.
\]

Hence,

\[
\boxed{E(AX^{(1)})=10}.
\]

\subsection*{(c) $\operatorname{Cov}(X^{(1)})$}

The covariance matrix for $X^{(1)}$ is the upper-left $2\times2$ block of $\Sigma_X$:

\[
\boxed{
\operatorname{Cov}(X^{(1)})
=
\begin{bmatrix}
3&0\\
0&1
\end{bmatrix}
}.
\]

\subsection*{(d) $\operatorname{Cov}(AX^{(1)})$}

Using

\[
\operatorname{Cov}(AX^{(1)})
=
A\operatorname{Cov}(X^{(1)})A',
\]

we obtain

\[
\operatorname{Cov}(AX^{(1)})
=
\begin{bmatrix}
1&2
\end{bmatrix}
\begin{bmatrix}
3&0\\
0&1
\end{bmatrix}
\begin{bmatrix}
1\\
2
\end{bmatrix}.
\]

First,

\[
\begin{bmatrix}
1&2
\end{bmatrix}
\begin{bmatrix}
3&0\\
0&1
\end{bmatrix}
=
\begin{bmatrix}
3&2
\end{bmatrix}.
\]

Thus,

\[
\operatorname{Cov}(AX^{(1)})
=
\begin{bmatrix}
3&2
\end{bmatrix}
\begin{bmatrix}
1\\
2
\end{bmatrix}
=
3+4=7.
\]

Therefore,

\[
\boxed{
\operatorname{Cov}(AX^{(1)})=7
}.
\]

\subsection*{(e) $E(X^{(2)})$}

The last two elements of the mean vector are

\[
\boxed{
E(X^{(2)})
=
\begin{bmatrix}
2\\
1
\end{bmatrix}
}.
\]

\subsection*{(f) $E(BX^{(2)})$}

Using

\[
E(BX^{(2)})=BE(X^{(2)}),
\]

we obtain

\[
E(BX^{(2)})
=
\begin{bmatrix}
1&-2\\
2&-1
\end{bmatrix}
\begin{bmatrix}
2\\
1
\end{bmatrix}.
\]

Therefore,

\[
E(BX^{(2)})
=
\begin{bmatrix}
2-2\\
4-1
\end{bmatrix}
=
\begin{bmatrix}
0\\
3
\end{bmatrix}.
\]

Hence,

\[
\boxed{
E(BX^{(2)})
=
\begin{bmatrix}
0\\
3
\end{bmatrix}
}.
\]

\subsection*{(g) $\operatorname{Cov}(X^{(2)})$}

The covariance matrix for $X^{(2)}$ is the lower-right $2\times2$ block of $\Sigma_X$:

\[
\boxed{
\operatorname{Cov}(X^{(2)})
=
\begin{bmatrix}
9&-2\\
-2&4
\end{bmatrix}
}.
\]

\subsection*{(h) $\operatorname{Cov}(BX^{(2)})$}

Using

\[
\operatorname{Cov}(BX^{(2)})
=
B\operatorname{Cov}(X^{(2)})B',
\]

we have

\[
\operatorname{Cov}(BX^{(2)})
=
\begin{bmatrix}
1&-2\\
2&-1
\end{bmatrix}
\begin{bmatrix}
9&-2\\
-2&4
\end{bmatrix}
\begin{bmatrix}
1&2\\
-2&-1
\end{bmatrix}.
\]

First,

\[
B\operatorname{Cov}(X^{(2)})
=
\begin{bmatrix}
1&-2\\
2&-1
\end{bmatrix}
\begin{bmatrix}
9&-2\\
-2&4
\end{bmatrix}
=
\begin{bmatrix}
13&-10\\
20&-8
\end{bmatrix}.
\]

Therefore,

\[
\operatorname{Cov}(BX^{(2)})
=
\begin{bmatrix}
13&-10\\
20&-8
\end{bmatrix}
\begin{bmatrix}
1&2\\
-2&-1
\end{bmatrix}.
\]

Thus,

\[
\operatorname{Cov}(BX^{(2)})
=
\begin{bmatrix}
13+20&26+10\\
20+16&40+8
\end{bmatrix}
=
\begin{bmatrix}
33&36\\
36&48
\end{bmatrix}.
\]

Hence,

\[
\boxed{
\operatorname{Cov}(BX^{(2)})
=
\begin{bmatrix}
33&36\\
36&48
\end{bmatrix}
}.
\]

\subsection*{(i) $\operatorname{Cov}(X^{(1)},X^{(2)})$}

The covariance between $X^{(1)}$ and $X^{(2)}$ is the upper-right $2\times2$ block of $\Sigma_X$:

\[
\boxed{
\operatorname{Cov}(X^{(1)},X^{(2)})
=
\begin{bmatrix}
2&2\\
1&0
\end{bmatrix}
}.
\]

\subsection*{(j) $\operatorname{Cov}(AX^{(1)},BX^{(2)})$}

Using

\[
\operatorname{Cov}(AX^{(1)},BX^{(2)})
=
A\operatorname{Cov}(X^{(1)},X^{(2)})B',
\]

we obtain

\[
\operatorname{Cov}(AX^{(1)},BX^{(2)})
=
\begin{bmatrix}
1&2
\end{bmatrix}
\begin{bmatrix}
2&2\\
1&0
\end{bmatrix}
\begin{bmatrix}
1&2\\
-2&-1
\end{bmatrix}.
\]

First,

\[
\begin{bmatrix}
1&2
\end{bmatrix}
\begin{bmatrix}
2&2\\
1&0
\end{bmatrix}
=
\begin{bmatrix}
4&2
\end{bmatrix}.
\]

Therefore,

\[
\operatorname{Cov}(AX^{(1)},BX^{(2)})
=
\begin{bmatrix}
4&2
\end{bmatrix}
\begin{bmatrix}
1&2\\
-2&-1
\end{bmatrix}.
\]

Thus,

\[
\operatorname{Cov}(AX^{(1)},BX^{(2)})
=
\begin{bmatrix}
4-4&8-2
\end{bmatrix}
=
\begin{bmatrix}
0&6
\end{bmatrix}.
\]

Hence,

\[
\boxed{
\operatorname{Cov}(AX^{(1)},BX^{(2)})
=
\begin{bmatrix}
0&6
\end{bmatrix}
}.
\]


\section*{Problem 2.41}

Given

\[
X=
\begin{bmatrix}
X_1\\
X_2\\
X_3\\
X_4
\end{bmatrix},
\qquad
\mu_X=
\begin{bmatrix}
3\\
2\\
-2\\
0
\end{bmatrix},
\]

and

\[
\Sigma_X=
\begin{bmatrix}
3&0&0&0\\
0&3&0&0\\
0&0&3&0\\
0&0&0&3
\end{bmatrix}.
\]

Also,

\[
A=
\begin{bmatrix}
1&-1&0&0\\
1&1&-2&0\\
1&1&1&-3
\end{bmatrix}.
\]

\subsection*{(a) Find $E(AX)$}

Using

\[
E(AX)=AE(X)=A\mu_X,
\]

we obtain

\[
E(AX)
=
\begin{bmatrix}
1&-1&0&0\\
1&1&-2&0\\
1&1&1&-3
\end{bmatrix}
\begin{bmatrix}
3\\
2\\
-2\\
0
\end{bmatrix}.
\]

Calculating each element,

\[
E(AX)
=
\begin{bmatrix}
1(3)-1(2)+0(-2)+0(0)\\
1(3)+1(2)-2(-2)+0(0)\\
1(3)+1(2)+1(-2)-3(0)
\end{bmatrix}.
\]

Therefore,

\[
E(AX)
=
\begin{bmatrix}
1\\
9\\
3
\end{bmatrix}.
\]

Hence,

\[
\boxed{
E(AX)=
\begin{bmatrix}
1\\
9\\
3
\end{bmatrix}
}.
\]

\subsection*{(b) Find $\operatorname{Cov}(AX)$}

For a linear transformation,

\[
\operatorname{Cov}(AX)
=
A\Sigma_XA'.
\]

Since

\[
\Sigma_X=3I_4,
\]

we have

\[
\operatorname{Cov}(AX)
=
3AA'.
\]

First,

\[
A'=
\begin{bmatrix}
1&1&1\\
-1&1&1\\
0&-2&1\\
0&0&-3
\end{bmatrix}.
\]

Therefore,

\[
AA'
=
\begin{bmatrix}
1&-1&0&0\\
1&1&-2&0\\
1&1&1&-3
\end{bmatrix}
\begin{bmatrix}
1&1&1\\
-1&1&1\\
0&-2&1\\
0&0&-3
\end{bmatrix}.
\]

Calculating the entries,

\[
AA'
=
\begin{bmatrix}
2&0&0\\
0&6&0\\
0&0&12
\end{bmatrix}.
\]

Thus,

\[
\operatorname{Cov}(AX)
=
3
\begin{bmatrix}
2&0&0\\
0&6&0\\
0&0&12
\end{bmatrix}
\]

and therefore

\[
\boxed{
\operatorname{Cov}(AX)
=
\begin{bmatrix}
6&0&0\\
0&18&0\\
0&0&36
\end{bmatrix}
}.
\]

The variances are

\[
\boxed{
\operatorname{Var}(Y_1)=6,\qquad
\operatorname{Var}(Y_2)=18,\qquad
\operatorname{Var}(Y_3)=36,
}
\]

where

\[
Y=AX.
\]

The covariances are

\[
\boxed{
\operatorname{Cov}(Y_1,Y_2)=0,
}
\]

\[
\boxed{
\operatorname{Cov}(Y_1,Y_3)=0,
}
\]

and

\[
\boxed{
\operatorname{Cov}(Y_2,Y_3)=0.
}
\]

\subsection*{(c) Which pairs of linear combinations have zero covariances?}

The three linear combinations are

\[
Y_1=X_1-X_2,
\]

\[
Y_2=X_1+X_2-2X_3,
\]

and

\[
Y_3=X_1+X_2+X_3-3X_4.
\]

From part (b),

\[
\operatorname{Cov}(AX)
=
\begin{bmatrix}
6&0&0\\
0&18&0\\
0&0&36
\end{bmatrix}.
\]

The off-diagonal entries are all zero. Therefore,

\[
\boxed{
\operatorname{Cov}(Y_1,Y_2)=0,
}
\]

\[
\boxed{
\operatorname{Cov}(Y_1,Y_3)=0,
}
\]

and

\[
\boxed{
\operatorname{Cov}(Y_2,Y_3)=0.
}
\]

Thus, every pair of the three linear combinations has zero covariance.

\end{document}
